% ============================================================
%  9. Conclusion
% ============================================================
\section{Conclusion}
\label{sec:conclusion}

We presented \LightRet{}, a lightweight, vocabulary-free named entity
recognition model trained via three-stage progressive knowledge distillation.
By grounding all word representations in the frozen pretrained \RetVec{}
character embedder, \LightRet{} eliminates vocabulary constraints entirely:
it processes arbitrary Unicode text without a tokenizer and is robust to the
character-level perturbations that cripple subword-based models.

The three-stage pipeline---sentence-level distillation into \RetBERT{},
token-level compression into a BiGRU--Transformer backbone, and
noisy-student NER fine-tuning with dynamic BIO label projection---transfers
BERT's contextual knowledge progressively into a model with
${\sim}4$M parameters, roughly $28{\times}$ smaller than BERT-base.

Our experiments on CoNLL-2003 reveal a clear trade-off that favours
\LightRet{} in real-world deployments: while BERT-base and DistilBERT
achieve strong clean-text F1 (91.25 and 89.93 respectively), they degrade
catastrophically under character noise---dropping to 52.1 and 43.0 F1 at
medium noise and below 25 F1 at high noise.
\LightRet{}, though modestly lower on clean text (85.7 F1), retains 78.0 F1
at medium noise and 69.2 F1 at high noise, striking the right balance
between clean-text competitiveness and noise robustness that neither
BERT-scale nor distilled subword models can match.
\LightRet{} thus substantially outperforms BERT-base and DistilBERT across
all noisy evaluation conditions.  The dynamic label projection
algorithm correctly handles word-splitting noise while preserving BIO validity,
and the compound distillation--classification loss in Stage~3 proves critical:
neither the classification signal nor the teacher alignment alone achieves
the same robustness.

We release all code, training scripts, pretrained weights, and the local
CoNLL-2003 and WikiText-103 dataset caches to facilitate reproducibility.
We hope \LightRet{} serves as a practical foundation for noise-robust NER in
production environments where BERT-scale models are cost-prohibitive.
